\documentclass[12pt, a4paper]{article}

% \usepackage[italian]{babel}  % requires texlive-lang-italian
\usepackage[utf8]{inputenc}
\usepackage[T1]{fontenc}
\usepackage{geometry}
\usepackage{hyperref}
\usepackage{booktabs}
\usepackage{array}
\usepackage{longtable}
\usepackage{graphicx}
\usepackage{xcolor}
\usepackage{enumitem}
\usepackage{parskip}
%\usepackage{microtype}
\usepackage{titlesec}
\usepackage{fancyhdr}
\usepackage{amsmath}

\geometry{
    top=2.5cm,
    bottom=2.5cm,
    left=3cm,
    right=3cm,
}

\hypersetup{
    colorlinks=true,
    linkcolor=blue!60!black,
    urlcolor=blue!60!black,
}

\titleformat{\section}{\large\bfseries}{}{0em}{\thesection.\quad}
\titleformat{\subsection}{\normalsize\bfseries}{}{0em}{\thesubsection.\quad}

\pagestyle{fancy}
\fancyhf{}
\rhead{\small Analisi Statistica -- RIPoSt}
\lhead{\small Disregolazione Emotiva in ADHD}
\cfoot{\thepage}

\newcommand{\rhosym}{\ensuremath{\varrho}}
\newcommand{\nota}[1]{\vspace{0.4em}\noindent\colorbox{gray!15}{\parbox{\dimexpr\linewidth-2\fboxsep}{\small\textbf{Nota clinica:} #1}}\vspace{0.4em}}

% ──────────────────────────────────────────────────────────────────────────────
\begin{document}

\begin{titlepage}
    \centering
    \vspace*{2cm}
    {\LARGE\bfseries Analisi Statistica della Disregolazione Emotiva\\[0.5em]
    in Pazienti Adulti con ADHD\par}
    \vspace{1.5cm}
    {\large Utilizzo del Questionario RIPoSt e della sua Versione di Screening (RIPoSt-SV)\par}
    \vspace{2cm}
    {\normalsize Documento tecnico a supporto della tesi di laurea\par}
    \vspace{0.5cm}
    {\normalsize Scritto per studenti di medicina -- non è richiesta una formazione statistica avanzata\par}
    \vspace{3cm}
    {\small Marzo 2026}
\end{titlepage}

\tableofcontents
\newpage

% ──────────────────────────────────────────────────────────────────────────────
\section{Introduzione e obiettivi dello studio}

\subsection{Contesto clinico}

La disregolazione emotiva (DE) è una delle caratteristiche più pervasive e invalidanti
dell'ADHD nell'adulto, spesso più impattante sulla qualità di vita dei sintomi attentivi o
iperattivi stessi. Si manifesta tipicamente come una combinazione di:

\begin{itemize}
    \item \textbf{Scatti d'ira e reazioni esplosive di breve durata} -- perdita transitoria del
          controllo, facile irritabilità, litigi frequenti
    \item \textbf{Labilità affettiva} -- oscillazioni rapide dell'umore tra normale, depresso o
          euforico, che durano ore o al massimo qualche giorno
    \item \textbf{Iper-reattività emotiva} -- incapacità di gestire lo stress ordinario, reazioni
          emotive sproporzionate che interferiscono con il problem solving
\end{itemize}

Nonostante la sua rilevanza clinica, la DE è stata storicamente difficile da misurare in
modo standardizzato, soprattutto in campioni psichiatrici eterogenei.

\subsection{Lo strumento RIPoSt}

Il \textbf{Questionario sulla Reattività, Intensità, Polarità e Stabilità (RIPoSt-40)} è uno
strumento self-report di 40 item, validato in italiano, che quantifica diversi aspetti della
DE. Gli item sono valutati su scala Likert da 1 (mai) a 6 (sempre) e si organizzano in
quattro sottoscale:

\begin{center}
\begin{tabular}{lp{9cm}}
\toprule
\textbf{Sottoscala} & \textbf{Cosa misura} \\
\midrule
Instabilità affettiva (AI) & Quanto l'umore oscilla nel tempo \\
Emotività negativa (N) & Tendenza a emozioni negative intense e frequenti \\
Emotività positiva (P) & Tendenza a emozioni positive intense e frequenti \\
Impulsività emotiva (EI) & Incapacità di controllare reazioni impulsive alle emozioni \\
\bottomrule
\end{tabular}
\end{center}

A partire da 30 dei 40 item (escludendo la sottoscala di emotività positiva), si calcola il
punteggio totale di \textbf{Disregolazione Emotiva Negativa (NED)}, con range 39--178.
Punteggi più alti indicano maggiore gravità della DE.

\subsection{La versione screening: RIPoSt-SV}

La \textbf{RIPoSt-SV} (Screening Version) è un algoritmo a sei item, derivato dalla RIPoSt-40
tramite un albero decisionale, che classifica ogni paziente come:

\begin{itemize}
    \item \textbf{CSED+ (True):} presenza di Disregolazione Emotiva Clinicamente Significativa
    \item \textbf{CSED- (False):} assenza di Disregolazione Emotiva Clinicamente Significativa
\end{itemize}

La CSED (Clinically Significant Emotional Dysregulation) è definita dalla presenza di almeno
2 dei 3 criteri qualitativi sopra descritti \textbf{E} una gravità di almeno 4/7 alla Clinical
Global Impression. Lo strumento è stato sviluppato e validato da Brancati et al. (2024,
\textit{Journal of Affective Disorders}) e ha mostrato un'accuratezza dell'83\% nel campione
di test.

\subsection{Obiettivi dell'analisi}

Il campione in studio è composto da \textbf{200 pazienti adulti con diagnosi di ADHD},
suddivisi nei sottotipi:

\begin{itemize}
    \item \textbf{INAT} -- sottotipo prevalentemente inattentivo (n = 98)
    \item \textbf{COMB} -- sottotipo combinato (n = 97)
    \item \textbf{IPER} -- sottotipo prevalentemente iperattivo (n = 4, solo descrittivo)
\end{itemize}

L'analisi si propone di rispondere a tre domande principali:

\begin{enumerate}
    \item Quali variabili cliniche e psicometriche sono \textit{associate} alla gravità della DE
          e allo stato CSED?
    \item Quali variabili sono \textit{predittori indipendenti} della DE, cioè aggiungono
          informazione unica anche quando si tengono in conto tutte le altre variabili?
    \item I pazienti si raggruppano spontaneamente in \textit{profili clinici} che coincidono
          con lo stato CSED? E questi profili predicono la CSED meglio o peggio dello strumento
          RIPoSt-SV?
\end{enumerate}

Ogni analisi viene eseguita su tre campioni: \textbf{campione completo} (n = 200),
\textbf{sottotipo INAT} e \textbf{sottotipo COMB}, per verificare se i risultati siano
consistenti o subtype-specifici.

% ──────────────────────────────────────────────────────────────────────────────
\section{Il dataset}

\subsection{Struttura generale}

Il dataset contiene 200 righe (pazienti) e 109 colonne (variabili). Ogni riga corrisponde a
un paziente anonimizzato. Le variabili si dividono in:

\begin{enumerate}
    \item \textbf{Variabili demografiche}
    \item \textbf{Variabili cliniche} (diagnosi, comorbilità, anamnesi familiare, farmaci)
    \item \textbf{Scale psicometriche} (punteggi continui di strumenti standardizzati)
    \item \textbf{Outcome} (RIPoSt-NED e RIPoSt-SV -- non utilizzati come predittori)
\end{enumerate}

\subsection{Variabili demografiche}

\begin{center}
\begin{tabular}{lp{9cm}}
\toprule
\textbf{Variabile} & \textbf{Descrizione} \\
\midrule
Genere & Sesso del paziente (M=1, F=0 nella codifica) \\
Age & Età in anni \\
school\_years & Anni di scolarità \\
repeated\_grades & Numero di anni scolastici ripetuti \\
marital\_status & Stato civile (celibe/nubile, sposato/a, divorziato/a) \\
\bottomrule
\end{tabular}
\end{center}

\nota{Il genere è stato codificato numericamente come M=1, F=0. Una correlazione negativa
con questa variabile significa che i maschi tendono ad avere punteggi più bassi -- ovvero le
femmine hanno DE più grave. Questo è coerente con la letteratura sull'ADHD.}

\subsection{Variabili cliniche}

Includono diagnosi di asse I secondo il DSM-5, storia familiare, uso di sostanze e
trattamento farmacologico. La maggior parte sono variabili \textbf{binarie} (vero/falso):

\begin{center}
\begin{tabular}{lp{8.5cm}}
\toprule
\textbf{Gruppo} & \textbf{Variabili} \\
\midrule
Storia familiare & FH\_bipolar, FH\_mood, FH\_neurodev, FH\_anxiety, FH\_psychosis, FH\_suicide, FH\_SUD \\
Disturbi dell'umore & mood (BD1/BD2/CYC/MDD/NO), mood\_any, mood\_onset \\
Uso di sostanze & alcool, opioids, cannabis, bdz, cocaine, stimulants, sud\_any \\
Disturbi d'ansia & panic, agoraphobia, social\_anx, GAD, OCD, anx\_any \\
Disturbi alimentari & anorexia, bulimia, BED, eat\_any \\
Altri & Tourette, tic, oppositional, conduct, conversion, somatic \\
Trattamento & Trattato (trattamento farmacologico in corso) \\
Anamnesi & first\_referral (età primo contatto psichiatrico), mood\_onset (età esordio umore) \\
\bottomrule
\end{tabular}
\end{center}

\subsection{Scale psicometriche}

Le scale sono punteggi \textbf{continui}: valori più alti indicano generalmente maggiore
gravità sintomatologica.

\begin{center}
\begin{longtable}{lp{5.5cm}p{4cm}}
\toprule
\textbf{Scala} & \textbf{Cosa misura} & \textbf{Colonne nel dataset} \\
\midrule
CAARS & Sintomi ADHD (self-report, Conners) & CAARS\_A/B/C/D \\
BRIEF & Difficoltà di funzionamento esecutivo & BRIEF\_GEC (totale), BRI, MI + 8 sottoscale \\
FAST & Disabilità funzionale globale & FAST\_TOT + 6 sottoscale \\
ASRS & Gravità ADHD (self-report adulti) & ASRS \\
LPFS & Funzionamento della personalità & LPFS\_S (sé), LPFS\_I (interpersonale) \\
HCL & Sintomi ipomaniacali & HCL\_TOT, HCL\_F1, HCL\_F2 \\
TEMPS & Temperamento affettivo & TEMPS\_D (depressivo), C (ciclotimico), H (ipertimico), I (irritabile), A (ansioso) \\
PID & Tratti di personalità (DSM-5) & PID\_NEG, DET, ANT, DIS, ANA, PSY \\
BIS & Impulsività (Barratt) & BIS\_TOT + BIS\_A/M/P \\
BRIAN & Ritmi biologici & BRIAN\_TOT + 5 sottoscale \\
\bottomrule
\end{longtable}
\end{center}

\subsection{Dati mancanti}

Non tutti i pazienti hanno completato tutte le scale. Le principali lacune sono:

\begin{itemize}
    \item CAARS: circa 39 valori mancanti ($\sim$20\%)
    \item BRIEF: circa 40 valori mancanti ($\sim$20\%)
    \item FAST\_TOT: circa 79 valori mancanti ($\sim$40\%)
\end{itemize}

\textbf{Strategia adottata: analisi per casi disponibili} (\textit{available-case analysis},
o \textit{pairwise deletion}). Ogni analisi usa tutti i pazienti che hanno dati validi per
quella specifica variabile, indipendentemente dai dati mancanti su altre variabili. In questo
modo non si perde una quota rilevante del campione. Per ogni risultato riportato, viene
indicato il numero di pazienti effettivamente inclusi (N).

\textbf{Non è stata eseguita alcuna imputazione dei dati mancanti.}

% ──────────────────────────────────────────────────────────────────────────────
\section{Analisi 1 -- Correlazione di Spearman con RIPoSt-NED}

\subsection{Domanda di ricerca}

\textit{Quali variabili sono associate alla gravità della disregolazione emotiva, misurata
dal punteggio continuo NED?}

\subsection{Principio del metodo}

La \textbf{correlazione di Spearman} misura quanto due variabili tendono a muoversi insieme,
in modo ordinato. Non assume che la relazione sia perfettamente lineare, né che i dati
seguano una distribuzione normale -- lavora sui ranghi, cioè sull'ordine relativo dei valori.
Per questo è particolarmente adatta a dati psichiatrici, che raramente sono distribuiti in
modo perfettamente normale.

Il risultato è il \textbf{coefficiente di correlazione $\varrho$ (rho)}, che varia tra
$-1$ e $+1$:

\begin{center}
\begin{tabular}{cc}
\toprule
\textbf{Valore di $|\varrho|$} & \textbf{Interpretazione} \\
\midrule
0.8 -- 1.0 & Molto forte \\
0.6 -- 0.8 & Forte \\
0.4 -- 0.6 & Moderata \\
0.2 -- 0.4 & Debole \\
0.0 -- 0.2 & Trascurabile \\
\bottomrule
\end{tabular}
\end{center}

Il \textbf{segno} indica la direzione:
\begin{itemize}
    \item $\varrho > 0$: al crescere di quella variabile, cresce anche NED (più sintomi = DE
          più grave)
    \item $\varrho < 0$: al crescere di quella variabile, NED diminuisce
\end{itemize}

\subsection{Esempio concreto}

Il temperamento ciclotimico (\textbf{TEMPS\_C}) ha $\varrho = +0.81$ nel campione completo.
Questo significa che i pazienti con punteggi più alti alla sottoscala del temperamento
ciclotimico hanno quasi sempre anche punteggi NED più elevati. La relazione è molto
consistente e forte.

Il punteggio \textbf{Genere} (M=1, F=0) ha $\varrho$ negativo: i maschi tendono ad avere
punteggi NED più bassi, ovvero le femmine hanno mediamente una DE più grave
(media NED femmine $\approx 127.6$ vs. maschi $\approx 108.9$).

La variabile \textbf{mood\_onset} (età di esordio del disturbo dell'umore) ha $\varrho$
negativo: esordio più precoce (numero più basso) è associato a DE più grave (NED più alto).

\subsection{Limiti di questa analisi}

Questa analisi guarda ogni variabile \textbf{una alla volta}. Non tiene conto del fatto che
molte variabili sono correlate tra loro. Per esempio, TEMPS\_C e ASRS sono entrambe
fortemente correlate con NED, ma sono anche correlate tra loro. Non possiamo sapere se
entrambe contribuiscono in modo indipendente, o se una è semplicemente un surrogato dell'altra.
Questo è il problema che risolve l'Analisi 2.

\subsection{Output generati}

Per ogni campione (completo, INAT, COMB) vengono generati:
\begin{itemize}
    \item \texttt{spearman\_results.csv} -- tabella completa con $\varrho$, p-value e N per
          ogni predittore
    \item \texttt{demographics.png} -- grafico a barre per le variabili demografiche
    \item \texttt{clinical.png} -- grafico a barre per le variabili cliniche
    \item \texttt{scale\_totals.png} -- grafico a barre per i totali delle scale psicometriche
    \item \texttt{scale\_subscales.png} -- grafico a barre per le sottoscale
\end{itemize}

Nei grafici, le barre \textbf{colorate} (blu scuro = correlazione positiva significativa,
rosso scuro = correlazione negativa significativa) indicano associazioni statisticamente
rilevanti (p $<$ 0.05); le barre chiare indicano associazioni non significative.

% ──────────────────────────────────────────────────────────────────────────────
\section{Analisi 2 -- Regressione Lineare LASSO su RIPoSt-NED}

\subsection{Domanda di ricerca}

\textit{Quali variabili predicono la gravità della DE in modo indipendente, cioè dopo aver
tenuto in conto tutte le altre variabili contemporaneamente?}

\subsection{Il problema della multicollinearità}

L'Analisi 1 mostra che molte variabili sono associate alla DE. Ma molte di queste variabili
sono anche associate tra loro. Per esempio, TEMPS\_C e ASRS sono entrambe elevate nei
pazienti con DE grave, ma sono anche correlate tra loro. La domanda quindi diventa:

\textit{Sapendo già il punteggio TEMPS\_C di un paziente, l'ASRS aggiunge qualcosa di nuovo?
O mi sta solo ridando la stessa informazione?}

La \textbf{regressione lineare multipla} risponde a questa domanda: include tutte le variabili
contemporaneamente e stima il contributo \textit{indipendente} di ciascuna.

\subsection{Il problema del sovraffittamento e la soluzione LASSO}

Con 98 predittori e solo 200 pazienti, la regressione standard è inaffidabile: il modello
troverebbe pattern casuali nei dati che non si riplicherebbero su nuovi pazienti. Questo
fenomeno si chiama \textbf{sovraffittamento} (\textit{overfitting}).

Il metodo \textbf{LASSO} (Least Absolute Shrinkage and Selection Operator) risolve questo
problema aggiungendo una \textit{penalità} per ogni variabile inclusa nel modello. Per
sopravvivere, una variabile deve contribuire abbastanza da giustificare il costo. Il
risultato: la maggior parte dei coefficienti vengono portati esattamente a zero, e solo le
variabili con un segnale indipendente reale sopravvivono.

\subsection{La cross-validazione a 5 fold}

La forza della penalità (parametro $\alpha$) non viene scelta manualmente, ma tramite
\textbf{cross-validazione a 5 fold}:

\begin{enumerate}
    \item I 200 pazienti vengono divisi in 5 gruppi da 40
    \item Il modello viene addestrato su 4 gruppi (160 pazienti) e testato sul 5° (40 pazienti)
    \item Questo si ripete 5 volte, lasciando ogni volta fuori un gruppo diverso
    \item Si sceglie il $\alpha$ che minimizza l'errore medio sui 5 test
\end{enumerate}

Questo garantisce che la performance riportata sia una stima onesta di come il modello si
comporterebbe su un paziente nuovo, mai visto prima.

\subsection{Cosa si legge dai risultati}

\begin{itemize}
    \item \textbf{Variabili sopravvissute al LASSO:} sono quelle con un contributo
          indipendente reale. Nel campione completo, 15 variabili su 98 sono sopravvissute.
    \item \textbf{Coefficiente $\beta$:} indica direzione e forza del contributo indipendente.
          $\beta > 0$ significa che punteggi più alti su quella variabile predicono NED più
          alto. $\beta < 0$ significa il contrario. Valori di $|\beta|$ più grandi indicano
          un contributo maggiore.
    \item \textbf{R² cross-validato $\approx 0.77$:} il modello spiega il 77\% della varianza
          nel punteggio NED. È un risultato molto buono per dati psichiatrici.
    \item \textbf{RMSE $\approx 14.5$:} in media, la previsione del modello si scosta di
          circa 14.5 punti dal valore reale (su una scala 39--178).
\end{itemize}

\subsection{Esempio concreto}

TEMPS\_C ha $\beta \approx +9.78$ e ASRS ha $\beta \approx +6.13$. Entrambe sopravvivono
al LASSO, il che significa che ciascuna contribuisce in modo indipendente alla gravità
della DE: sapere il temperamento ciclotimico di un paziente non rende inutile conoscere
la sua gravità ADHD all'ASRS, e viceversa. Sono due informazioni distinte.

Al contrario, molte variabili che correlano con NED nell'Analisi 1 scompaiono nell'Analisi 2,
perché la loro associazione con NED è mediata da TEMPS\_C o ASRS: non aggiungono nulla di
nuovo.

\nota{Una variabile eliminata dal LASSO non è necessariamente priva di interesse clinico.
Potrebbe semplicemente veicolare la stessa informazione di una variabile sopravvissuta.
Per esempio, se TEMPS\_C e HCL\_TOT (sintomi ipomaniacali) sono molto correlati tra loro,
LASSO tende a tenere solo uno dei due. Questo non significa che HCL\_TOT sia clinicamente
irrilevante.}

\subsection{Output generati}

Per ogni campione:
\begin{itemize}
    \item \texttt{lasso\_coefficients.csv} -- variabili sopravvissute con i relativi $\beta$
    \item \texttt{lasso\_metrics.csv} -- R², RMSE, parametro $\alpha$ ottimale, N
    \item \texttt{lasso\_coefficients.png} -- grafico a barre dei $\beta$ sopravvissuti
\end{itemize}

% ──────────────────────────────────────────────────────────────────────────────
\section{Analisi 3 -- Analisi di Associazione: CSED+ vs CSED-}

\subsection{Domanda di ricerca}

\textit{Per ogni variabile, i pazienti che risultano positivi allo screening CSED (RIPoSt-SV
= True) sono diversi da quelli negativi?}

\subsection{Logica del confronto tra gruppi}

Le Analisi 1 e 2 usavano NED come outcome \textbf{continuo} (gravità). L'Analisi 3 cambia
prospettiva e usa la \textbf{RIPoSt-SV} come outcome \textbf{binario}: ogni paziente è
CSED+ o CSED-.

I 200 pazienti vengono divisi in due gruppi: CSED+ (n = 132) e CSED- (n = 68). Per ogni
variabile predittrice, si testa se i due gruppi differiscono significativamente.

\subsection{Due test diversi per due tipi di variabili}

\subsubsection{Variabili continue (punteggi delle scale)}

Test utilizzato: \textbf{Mann-Whitney U}

Questo test non assume che i dati seguano una distribuzione normale. Confronta i ranghi
dei valori nei due gruppi: se i pazienti CSED+ tendono ad avere ranghi sistematicamente
più alti su quella scala, il test risulta significativo.

Misura dell'effetto: \textbf{correlazione rank-biseriale r}, che varia tra $-1$ e $+1$:
\begin{itemize}
    \item $r > 0$: i pazienti CSED+ hanno \textbf{punteggi più alti} su quella variabile
    \item $r < 0$: i pazienti CSED+ hanno \textbf{punteggi più bassi} su quella variabile
    \item $|r| \geq 0.5$: effetto grande
\end{itemize}

\textbf{Esempio:} TEMPS\_C ha $r = +0.64$. I pazienti CSED+ hanno punteggi di temperamento
ciclotimico significativamente più alti di quelli CSED-.

\textbf{Esempio:} mood\_onset (età di esordio del disturbo dell'umore) ha $r$ negativo.
I pazienti CSED+ hanno avuto un esordio più precoce del disturbo dell'umore -- l'esordio
precoce è associato alla DE clinicamente significativa.

\subsubsection{Variabili binarie/categoriali (comorbilità, storia familiare)}

Test utilizzato: \textbf{Chi-quadro} o \textbf{test esatto di Fisher} (quando i numeri
attesi sono piccoli)

Questi test verificano se la frequenza di una condizione (per esempio: presenza di un
disturbo d'ansia) è distribuita diversamente tra CSED+ e CSED-.

Misura dell'effetto: \textbf{V di Cramér}, che varia tra 0 e 1:
\begin{itemize}
    \item 0: nessuna associazione
    \item 1: associazione perfetta
    \item \textbf{Attenzione:} la V di Cramér non ha segno -- indica la \textit{forza}
          dell'associazione ma non la direzione. Per sapere se una condizione è più frequente
          in CSED+ o CSED-, occorre guardare le frequenze nel file CSV.
\end{itemize}

\textbf{Esempio:} mood\_any (presenza di qualsiasi disturbo dell'umore) ha una V di Cramér
significativa. Questo significa che la distribuzione dei disturbi dell'umore è diversa tra
i due gruppi. Dalla tabella delle frequenze si evince in quale gruppo siano più rappresentati.

\subsection{Limiti di questa analisi}

Come l'Analisi 1, anche questa analisi testa ogni variabile \textbf{individualmente}. Una
variabile può risultare significativa semplicemente perché è correlata con un'altra variabile
che è il vero driver. L'Analisi 4 affronta questo problema.

\subsection{Output generati}

Per ogni campione:
\begin{itemize}
    \item \texttt{association\_results.csv} -- tabella completa con dimensione d'effetto,
          p-value, N e test utilizzato per ogni predittore
    \item \texttt{demographics.png}, \texttt{clinical.png}, \texttt{scale\_totals.png},
          \texttt{scale\_subscales.png} -- grafici a barre per gruppo di variabili
\end{itemize}

Nei grafici, le barre colorate indicano significatività statistica (p $<$ 0.05):
blu scuro = effetto positivo significativo; rosso scuro = effetto negativo significativo.

% ──────────────────────────────────────────────────────────────────────────────
\section{Analisi 4 -- Regressione Logistica LASSO su RIPoSt-SV}

\subsection{Domanda di ricerca}

\textit{Quali variabili predicono in modo indipendente l'essere CSED+, e di quanto aumentano
o diminuiscono le probabilità di un esito positivo allo screening?}

\subsection{Differenza rispetto all'Analisi 3}

L'Analisi 3 guardava ogni variabile \textbf{una alla volta}. L'Analisi 4 le include
\textbf{tutte insieme} in un unico modello, permettendo di identificare quelle che contribuiscono
in modo genuinamente indipendente.

\textbf{Esempio pratico:} TEMPS\_C e ASRS sono entrambe più alte nei CSED+ (Analisi 3).
Ma sono anche correlate tra loro. L'Analisi 4 risponde: \textit{"Sapendo già TEMPS\_C,
l'ASRS cambia qualcosa alla nostra previsione di essere CSED+?"} Se sì, entrambe sopravvivono
al modello. Se no, una delle due viene eliminata.

\subsection{Principio della regressione logistica}

La regressione logistica è il metodo standard per predire un outcome binario (sì/no,
positivo/negativo) a partire da più variabili. Produce un modello che stima, per ogni
paziente, la \textbf{probabilità di essere CSED+} in base al suo profilo clinico e
psicometrico. Come nell'Analisi 2, viene applicata la penalità LASSO per eliminare le
variabili ridondanti.

\subsection{Come si interpretano i risultati}

Il risultato chiave è l'\textbf{odds ratio (OR)}:

\begin{equation*}
\text{OR} = e^{\beta}
\end{equation*}

\begin{itemize}
    \item \textbf{OR $>$ 1:} quella variabile \textbf{aumenta} le probabilità di essere CSED+
    \item \textbf{OR $<$ 1:} quella variabile \textbf{diminuisce} le probabilità di essere CSED+
    \item \textbf{OR = 1:} nessun effetto
\end{itemize}

\textbf{Esempio:} Nel campione completo, TEMPS\_C ha OR $\approx$ 1.76. Questo significa
che, per ogni punto in più nel punteggio del temperamento ciclotimico, le probabilità di
risultare CSED+ sono 1.76 volte più alte -- indipendentemente da ASRS, BRIAN\_ACT e tutte
le altre variabili nel modello.

\textbf{Attenzione all'interpretazione:} gli OR qui riportati si riferiscono a variabili
\textit{standardizzate}, quindi l'OR non è per un incremento di una singola unità grezza,
ma per un incremento di una deviazione standard. Questo permette di confrontare l'entità
degli effetti tra scale con unità diverse.

\subsection{Performance del modello}

\begin{itemize}
    \item \textbf{AUC cross-validata $\approx$ 0.80:} l'AUC (Area Under the Curve) misura la
          capacità del modello di distinguere CSED+ da CSED-. 0.5 = casuale, 1.0 = perfetto.
          0.80 indica una buona discriminazione.
    \item \textbf{Accuratezza cross-validata $\approx$ 76\%:} il modello classifica
          correttamente circa 3 pazienti su 4. Per confronto, la RIPoSt-SV raggiunge l'83\%
          nel campione di test della pubblicazione originale.
\end{itemize}

\nota{Il fatto che il modello logistico (76\%) sia meno accurato della RIPoSt-SV (83\%) è
un risultato clinicamente rilevante, non un fallimento dell'analisi. Significa che il profilo
clinico e psicometrico complessivo non è in grado di identificare la CSED con la stessa
precisione dello strumento specificamente costruito per farlo. Questo \textbf{valida la
RIPoSt-SV} come strumento che cattura qualcosa di specifico, non riducibile al profilo
clinico generale.}

\subsection{Output generati}

Per ogni campione:
\begin{itemize}
    \item \texttt{logistic\_coefficients.csv} -- variabili sopravvissute con $\beta$ e OR
    \item \texttt{logistic\_metrics.csv} -- AUC, accuratezza, C ottimale, N, CSED+/CSED-
    \item \texttt{logistic\_coefficients.png} -- grafico a barre con $\beta$ e OR indicati
          nelle etichette
\end{itemize}

% ──────────────────────────────────────────────────────────────────────────────
\section{Analisi 5 -- Clustering e Confronto con RIPoSt-SV}

\subsection{Domanda di ricerca}

\textit{I pazienti con ADHD si raggruppano spontaneamente in profili clinici distinti? E
questi profili predicono la CSED meglio o peggio della RIPoSt-SV?}

\subsection{Una logica completamente diversa}

Tutte le analisi precedenti partivano da un outcome (NED o RIPoSt-SV) e chiedevano quali
variabili lo predicessero. Il clustering è \textbf{non supervisionato}: non guarda l'outcome.
Guarda solo i pazienti e chiede:

\textit{"Guardando esclusivamente i punteggi delle scale psicometriche, questi 200 pazienti
si assomigliano tutti, o si dividono in gruppi naturali con profili distinti?"}

Solo \textit{dopo} aver identificato i gruppi, si va a vedere quanti CSED+ ci sono in
ciascun gruppo.

\subsection{Come funziona il K-means}

L'algoritmo \textbf{K-means} divide i pazienti in k gruppi (cluster) in modo che:
\begin{itemize}
    \item I pazienti \textbf{nello stesso cluster} siano il più simili possibile tra loro
          (punteggi simili su tutte le scale)
    \item I pazienti in \textbf{cluster diversi} siano il più diversi possibile
\end{itemize}

Il numero di cluster (k) non viene scelto a priori, ma viene determinato automaticamente
testando da k=2 a k=6 e scegliendo il valore che produce i gruppi più netti e ben separati,
misurato dal \textbf{punteggio silhouette} (più alto = gruppi meglio definiti). In tutti i
campioni analizzati, k=2 è risultato ottimale.

\subsection{Cosa si fa dopo il clustering}

Una volta assegnati i pazienti ai cluster:

\begin{enumerate}
    \item Si calcola la \textbf{percentuale di CSED+} in ogni cluster
    \item Si assegna a ogni cluster l'\textbf{etichetta di maggioranza} (se il 70\% dei
          pazienti in quel cluster è CSED+, il cluster viene etichettato "CSED+")
    \item Si usa questa etichetta come classificatore e si calcola l'\textbf{accuratezza} --
          quanti pazienti vengono classificati correttamente rispetto alla loro vera RIPoSt-SV
    \item Si confronta questa accuratezza con quella di riferimento della RIPoSt-SV (83\%)
\end{enumerate}

\subsection{I risultati e la loro interpretazione}

L'accuratezza del clustering è risultata del \textbf{66--69\%} in tutti i campioni,
significativamente inferiore all'83\% della RIPoSt-SV.

Questo è un risultato clinicamente importante. Significa che \textbf{il raggruppamento
naturale dei pazienti ADHD in base al loro profilo clinico e psicometrico non coincide con
la distinzione tra CSED+ e CSED-}. Essere CSED+ non equivale semplicemente all'essere un
paziente ADHD ``più grave'' su tutte le dimensioni: è una caratteristica clinica specifica
che richiede uno strumento dedicato per essere identificata in modo affidabile.

Questo risultato \textbf{valida la RIPoSt-SV}: lo strumento cattura qualcosa di specifico
che la profilazione clinica generale non riesce a cogliere con la stessa precisione.

\subsection{Il profilo dei cluster}

L'aspetto più interessante dal punto di vista clinico è l'\textbf{heatmap del profilo dei
cluster}: mostra per ogni scala se i pazienti di quel cluster punteggiano sopra o sotto la
media (standardizzata). Permette di rispondere a domande come:
\begin{itemize}
    \item Quali scale distinguono meglio i due cluster?
    \item Il cluster con più CSED+ ha anche impulsività più alta? Ritmi biologici più alterati?
          Temperamento più ciclotimico?
\end{itemize}

\subsection{Output generati}

Per ogni campione:
\begin{itemize}
    \item \texttt{silhouette\_scores.png} -- punteggio silhouette per k da 2 a 6 (motiva la
          scelta di k)
    \item \texttt{cluster\_profiles.png} -- heatmap dei punteggi medi standardizzati per cluster
    \item \texttt{csed\_rate\_per\_cluster.png} -- percentuale CSED+ per cluster
    \item \texttt{cluster\_assignments.csv} -- cluster e etichetta predetta per ogni paziente
    \item \texttt{clustering\_metrics.csv} -- k ottimale, silhouette, accuratezza, confronto 83\%
    \item \texttt{confusion\_matrix.csv} -- etichette vere vs. predette (CSED+ / CSED-)
\end{itemize}

% ──────────────────────────────────────────────────────────────────────────────
\section{Riepilogo delle analisi}

\begin{center}
\begin{longtable}{p{1.5cm}p{3cm}p{3.5cm}p{5cm}}
\toprule
\textbf{Analisi} & \textbf{Outcome} & \textbf{Metodo} & \textbf{Domanda a cui risponde} \\
\midrule
1 & RIPoSt-NED (continuo) & Correlazione di Spearman & Quali variabili variano insieme alla gravità della DE? \\
\addlinespace
2 & RIPoSt-NED (continuo) & Regressione lineare LASSO & Quali variabili predicono la gravità in modo indipendente? \\
\addlinespace
3 & RIPoSt-SV (binario) & Mann-Whitney / Chi-quadro + dimensione d'effetto & CSED+ e CSED- differiscono su questa variabile? \\
\addlinespace
4 & RIPoSt-SV (binario) & Regressione logistica LASSO & Quali variabili predicono indipendentemente la CSED? \\
\addlinespace
5 & Nessuno (non supervisionato) & K-means + confronto & I pazienti si raggruppano in profili che coincidono con CSED? \\
\bottomrule
\end{longtable}
\end{center}

Ogni analisi viene eseguita su tre campioni:

\begin{itemize}
    \item \textbf{Campione completo} (n = 200): massima potenza statistica, risultato di riferimento
    \item \textbf{Sottotipo INAT} (n = 98): risultati specifici per pazienti inattentivi
    \item \textbf{Sottotipo COMB} (n = 97): risultati specifici per pazienti con sottotipo combinato
\end{itemize}

% ──────────────────────────────────────────────────────────────────────────────
\section{Note metodologiche finali}

\subsection{Significatività statistica e dimensione dell'effetto}

Un risultato statisticamente significativo (p $<$ 0.05) non garantisce rilevanza clinica.
Con campioni di 200 pazienti, anche effetti molto piccoli possono risultare significativi.
Per questo motivo, in tutte le analisi vengono riportate le \textbf{dimensioni dell'effetto}
(correlazione di Spearman, rank-biserial r, V di Cramér, OR) che permettono di valutare la
\textit{grandezza pratica} della differenza, indipendentemente dalla significatività.

\subsection{Analisi univariate vs. multivariate}

Le Analisi 1 e 3 sono \textbf{univariate}: testano ogni variabile in isolamento. Un risultato
significativo non implica causalità o indipendenza dal resto del profilo clinico.

Le Analisi 2 e 4 sono \textbf{multivariate}: controllano per tutte le variabili simultaneamente.
Sono più conservative ma più informative sulla struttura causale.

La lettura integrata di analisi univariate e multivariate è la più ricca: le variabili che
risultano significative sia in isolamento (Analisi 1/3) che nel modello multivariato
(Analisi 2/4) sono quelle con il segnale più robusto e indipendente.

\subsection{Confronto tra sottotipi}

Variabili significative nel campione completo ma non nei sottotipi separati possono riflettere
semplicemente la riduzione di potenza statistica (da n=200 a n$\approx$98). Variabili
significative in un sottotipo ma non nell'altro possono invece indicare genuina eterogeneità
biologica o clinica tra INAT e COMB.

\bigskip

\hrule
\bigskip

\noindent\small\textbf{Riferimento bibliografico:}

Brancati GE, De Rosa U, Acierno D, et al. (2024). Development of a self-report screening
instrument for emotional dysregulation: the Reactivity, Intensity, Polarity and Stability
questionnaire, screening version (RIPoSt-SV). \textit{Journal of Affective Disorders},
355, 406--414. \url{https://doi.org/10.1016/j.jad.2024.03.167}

\end{document}
